% GENERATED FILE. Edit canonical lessons, then run npm run generate:books.
% canonical-source-hash: fnv1a64:f128c608995cfdae
% canonical-lessons: MW-C199-mohar, MW-C199-daskhat, MW-C199-rasid, MW-C199-lekho, MW-C199-madno

\chapter{Seal and Signature}
\label{ch:mw-seal-and-signature}
\hlchaptermodality{199}

\begin{chapteropening}
\textbf{By the end of this chapter:} \emph{I can say what a form still needs.}

Complete the last lesson of chapter 199: I can say what a form still needs.
\end{chapteropening}

\section[\texorpdfstring{\mw{मोहर} (mohar)}{मोहर (mohar)}]{\mw{मोहर} (mohar) — an official stamp; a seal}

\label{lesson:MW-C199-mohar}



\begin{quote}
Before the new one: say the Marwadi for a request, then the Marwadi for to explain.
\end{quote}

\subsection*{You'll want to know: \mw{मोहर}}
\textbf{\mw{मोहर}} — \emph{mohar} — \textquotedblleft{}an official stamp; a seal\textquotedblright{}.

\begin{quote}
\mw{अरजी} \mw{माथै} \mw{मोहर} \mw{कोनी।}
\end{quote}

\begin{quote}
\emph{arjī māthai mohar konī.}
\end{quote}

\begin{quote}
There is no stamp on the application.
\end{quote}

A feminine noun.

\begin{cousinweb}
\textbf{\mw{मोहर}} (\emph{mohar}) is usually traced to Persian \emph{muhr}, \textquotedblleft{}seal, stamp\textquotedblright{}.
\end{cousinweb}

\begin{grammarlens}[title={What you've built}]
One more word for reading notices and messages.
\end{grammarlens}

\subsection*{Your turn}
\begin{itemize}
\raggedright
  \item \emph{Say it:} \emph{mohar}
  \item \emph{Say it:} \emph{mohar}, once more
  \item \emph{Say it:} \emph{arjī māthai mohar konī.} — There is no stamp on the application.
  \item \emph{Recall:} say the Marwadi for a request, then the Marwadi for to explain, then say \emph{mohar} again
\end{itemize}

\begin{tcolorbox}[breakable,colback=teal!4,colframe=teal!35!black,title={Before you move on}]
What does \mw{मोहर} mean? (\textquotedblleft{}An official stamp; a seal\textquotedblright{}.) Say it once more.
\end{tcolorbox}
\section[\texorpdfstring{\mw{दसखत} (daskhat)}{दसखत (daskhat)}]{\mw{दसखत} (daskhat) — a signature}

\label{lesson:MW-C199-daskhat}



\begin{quote}
Before the new one: say the Marwadi for to explain, then the Marwadi for an official stamp.
\end{quote}

\subsection*{You'll want to know: \mw{दसखत}}
\textbf{\mw{दसखत}} — \emph{daskhat} — \textquotedblleft{}a signature\textquotedblright{}.

\begin{quote}
\mw{अठै} \mw{दसखत} \mw{करो।}
\end{quote}

\begin{quote}
\emph{aṭhai daskhat karo.}
\end{quote}

\begin{quote}
Sign here, please.
\end{quote}

A masculine noun.

\begin{cousinweb}
\textbf{\mw{दसखत}} (\emph{daskhat}) is usually traced to Persian \emph{dastkhaṭṭ}, \textquotedblleft{}handwriting\textquotedblright{}.
\end{cousinweb}

\begin{grammarlens}[title={What you've built}]
One more word for reading notices and messages.
\end{grammarlens}

\subsection*{Your turn}
\begin{itemize}
\raggedright
  \item \emph{Say it:} \emph{daskhat}
  \item \emph{Say it:} \emph{daskhat}, once more
  \item \emph{Say it:} \emph{aṭhai daskhat karo.} — Sign here, please.
  \item \emph{Recall:} say the Marwadi for to explain, then the Marwadi for an official stamp, then say \emph{daskhat} again
\end{itemize}

\begin{tcolorbox}[breakable,colback=teal!4,colframe=teal!35!black,title={Before you move on}]
What does \mw{दसखत} mean? (\textquotedblleft{}A signature\textquotedblright{}.) Say it once more.
\end{tcolorbox}
\section[\texorpdfstring{\mw{रसीद} (rasīd)}{रसीद (rasīd)}]{\mw{रसीद} (rasīd) — a receipt}

\label{lesson:MW-C199-rasid}



\begin{quote}
Before the new one: say the Marwadi for an official stamp, then the Marwadi for a signature.
\end{quote}

\subsection*{You'll want to know: \mw{रसीद}}
\textbf{\mw{रसीद}} — \emph{rasīd} — \textquotedblleft{}a receipt\textquotedblright{}.

\begin{quote}
\mw{आप} \mw{रसीद} \mw{राखो।}
\end{quote}

\begin{quote}
\emph{āp rasīd rākho.}
\end{quote}

\begin{quote}
Please keep the receipt.
\end{quote}

A feminine noun.

\begin{cousinweb}
\textbf{\mw{रसीद}} (\emph{rasīd}) is usually traced to Persian \emph{rasīd}, \textquotedblleft{}received; a receipt\textquotedblright{}.
\end{cousinweb}

\begin{grammarlens}[title={What you've built}]
One more word for reading notices and messages.
\end{grammarlens}

\subsection*{Your turn}
\begin{itemize}
\raggedright
  \item \emph{Say it:} \emph{rasīd}
  \item \emph{Say it:} \emph{rasīd}, once more
  \item \emph{Say it:} \emph{āp rasīd rākho.} — Please keep the receipt.
  \item \emph{Recall:} say the Marwadi for an official stamp, then the Marwadi for a signature, then say \emph{rasīd} again
\end{itemize}

\begin{tcolorbox}[breakable,colback=teal!4,colframe=teal!35!black,title={Before you move on}]
What does \mw{रसीद} mean? (\textquotedblleft{}A receipt\textquotedblright{}.) Say it once more.
\end{tcolorbox}
\section[\texorpdfstring{\mw{लेखो} (lekho)}{लेखो (lekho)}]{\mw{लेखो} (lekho) — an account; a reckoning}

\label{lesson:MW-C199-lekho}



\begin{quote}
Before the new one: say the Marwadi for a signature, then the Marwadi for a receipt.
\end{quote}

\subsection*{You'll want to know: \mw{लेखो}}
\textbf{\mw{लेखो}} — \emph{lekho} — \textquotedblleft{}an account; a reckoning\textquotedblright{}.

\begin{quote}
\mw{म्हारो} \mw{लेखो} \mw{पूरो} \mw{है।}
\end{quote}

\begin{quote}
\emph{mhāro lekho pūro hai.}
\end{quote}

\begin{quote}
My account is settled.
\end{quote}

A masculine noun.

\begin{cousinweb}
\textbf{\mw{लेखो}} (\emph{lekho}) is inherited from Sanskrit \emph{lekhya}, \textquotedblleft{}a writing, a document\textquotedblright{}; CDIAL lists Old Marwari \emph{lekha}, \textquotedblleft{}writings\textquotedblright{}, and Hindi \emph{lekhā}, \textquotedblleft{}account\textquotedblright{} (CDIAL 11108).
\end{cousinweb}

\begin{grammarlens}[title={What you've built}]
One more word for reading notices and messages.
\end{grammarlens}

\subsection*{Your turn}
\begin{itemize}
\raggedright
  \item \emph{Say it:} \emph{lekho}
  \item \emph{Say it:} \emph{lekho}, once more
  \item \emph{Say it:} \emph{mhāro lekho pūro hai.} — My account is settled.
  \item \emph{Recall:} say the Marwadi for a signature, then the Marwadi for a receipt, then say \emph{lekho} again
\end{itemize}

\begin{tcolorbox}[breakable,colback=teal!4,colframe=teal!35!black,title={Before you move on}]
What does \mw{लेखो} mean? (\textquotedblleft{}An account; a reckoning\textquotedblright{}.) Say it once more.
\end{tcolorbox}
\section[\texorpdfstring{\mw{मांडणो} (mā̃ḍṇo)}{मांडणो (mā̃ḍṇo)}]{\mw{मांडणो} (mā̃ḍṇo) — to write down; to enter (in a ledger)}

\label{lesson:MW-C199-madno}



\begin{quote}
Before the new one: say the Marwadi for a receipt, then the Marwadi for a reckoning.
\end{quote}

\subsection*{You'll want to know: \mw{मांडणो}}
\textbf{\mw{मांडणो}} — \emph{mā̃ḍṇo} — \textquotedblleft{}to write down; to enter (in a ledger)\textquotedblright{}.

\begin{quote}
\mw{मास्टर} \mw{नाम} \mw{मांडै} \mw{है।}
\end{quote}

\begin{quote}
\emph{māsṭar nām mā̃ḍai hai.}
\end{quote}

\begin{quote}
The teacher writes the names down.
\end{quote}

Past (I …): \mw{म्हैं} \mw{नाम} \mw{मांड्यो}, \emph{mhaĩ nām mā̃ḍyo}.

\begin{cousinweb}
\textbf{\mw{मांडणो}} (\emph{mā̃ḍṇo}) is inherited from Sanskrit \emph{maṇḍayati}, \textquotedblleft{}adorns, decorates\textquotedblright{}; CDIAL lists an Old Marwari form of it glossed \textquotedblleft{}writes\textquotedblright{} (CDIAL 9741).
\end{cousinweb}

\begin{grammarlens}[title={What you've built}]
One more word for reading notices and messages.
\end{grammarlens}

\subsection*{Your turn}
\begin{itemize}
\raggedright
  \item \emph{Say it:} \emph{mā̃ḍṇo}
  \item \emph{Say it:} \emph{mā̃ḍṇo}, once more
  \item \emph{Say it:} \emph{māsṭar nām mā̃ḍai hai.} — The teacher writes the names down.
  \item \emph{Recall:} say the Marwadi for a receipt, then the Marwadi for a reckoning, then say \emph{mā̃ḍṇo} again
\end{itemize}

\begin{tcolorbox}[breakable,colback=teal!4,colframe=teal!35!black,title={Before you move on}]
What does \mw{मांडणो} mean? (\textquotedblleft{}To write down; to enter (in a ledger)\textquotedblright{}.) Say it once more.
\end{tcolorbox}
